\setcounter{section}{1}

  \zwischenueberschrift{Hipermuka}

Misalkan
\maabbdisp {h} {\R^n } { \R
} {}
adalah sebuah
\definitionsverweis {fungsi terdiferensialkan}{}{}
dan

\mavergleichskette
{\vergleichskette
{ Y
}
{ = }{ h^{-1}(c)
}
{  }{
}
{    }{
}
{   }{
}
}
{}{}{}
adalah serat di atas

\mavergleichskette
{\vergleichskette
{ c
}
{ \in }{ \R
}
{  }{
}
{    }{
}
{   }{
}
}
{}{}{.}
Himpunan-himpunan bagian ini dipelajari khususnya dalam konteks teorema fungsi implisit. Jika
\definitionsverweis {diferensial total}{}{}
\maabbdisp {\left(D h \right)_{ P}} {\R^n } { \R
} {}
surjektif untuk setiap titik

\mavergleichskette
{\vergleichskette
{ P
}
{ \in }{ Y
}
{  }{
}
{    }{
}
{   }{
}
}
{}{}{}
---dengan kata lain, jika $P$ selalu merupakan
\definitionsverweis {titik reguler}{}{}
dari pemetaan itu---maka $Y$, dalam suatu lingkungan terbuka dari $P$, homeomorfik dengan suatu himpunan terbuka di $\R^{n-1}$. Selain itu, dalam konteks ini kita telah memperkenalkan
\definitionsverweis {ruang tangen pada serat}{}{}
sebagai kernel diferensial total. Ruang ini merupakan subruang vektor berdimensi $(n-1 )$ dari $\R^n$, tetapi biasanya dibayangkan melekat pada titik $P$, yakni sebagai ruang afin-linear. Di titik $P$, ruang tangen ini menyinggung $Y$. Kita menyebut $Y$ sebuah \stichwort {hipermuka} {,} karena ruang tangennya berdimensi satu lebih rendah daripada ruang sekitarnya. Untuk

\mavergleichskette
{\vergleichskette
{ n
}
{ = }{ 3
}
{  }{
}
{    }{
}
{   }{
}
}
{}{}{}
memang diperoleh sebuah permukaan. Jika kita menggunakan hasil kali dalam standar pada $\R^n$, ruang tangen pada serat juga dapat dipandang sebagai ruang semua vektor yang tegak lurus terhadap
\definitionsverweis {gradien}{}{}

\mathl{\operatorname{Grad} \,  h  (  P )}{}.

\inputbeispiel{}
{

Misalkan

\mavergleichskettedisp
{\vergleichskette
{ h(x,y,z)
}
{ =} { x^2+y^2+z^2
}
{ } {
}
{ } {
}
{ } {
}
}
{}{}{}
dan kita tinjau serat $Y$ dari $h$ di atas $1$, yaitu permukaan bola berjari-jari $1$ dengan pusat di titik asal. Diferensial totalnya adalah
\mathl{\left(  2x  , \,   2y  , \,   2z                 \right)}{,} sehingga di setiap titik permukaan bola setidaknya satu komponennya tidak sama dengan $0$; oleh karena itu, $h$
\definitionsverweis {reguler}{}{}
di setiap titik pada $Y$. Misalkan

\mavergleichskette
{\vergleichskette
{ P
}
{ = }{ \left(  a   , \,   b  , \,  c                 \right)
}
{ \in }{ Y
}
{    }{
}
{   }{
}
}
{}{}{.}
Ruang tangen di $P$ diberikan oleh syarat linear

\mavergleichskettedisp
{\vergleichskette
{ ax+by+cz
}
{ =} { 0
}
{ } {
}
{ } {
}
{ } {
}
}
{}{}{}
dan merupakan himpunan semua vektor yang ortogonal terhadap gradien
\mathl{\begin{pmatrix}  2a \\ 2b\\  2c   \end{pmatrix}}{}.

}

\inputbeispiel{}
{

$\,$

\bild{ \begin{center}
\includegraphics[width=5.5cm]{ \bildeinlesungsvg {3d-function-6} {svg} }
\end{center}
\bildtext {} }

\bildlizenz { 3d-function-6.svg } {} {MartinThoma} {Commons} {CC BY 3.0} {}

Misalkan

\mavergleichskette
{\vergleichskette
{ V
}
{ \subseteq }{ \R^{n-1}
}
{  }{
}
{    }{
}
{   }{
}
}
{}{}{}
adalah sebuah
\definitionsverweis {himpunan terbuka}{}{}
dan misalkan
\maabbdisp {f} { V } { \R
} {}
adalah sebuah
\definitionsverweis {fungsi yang terdiferensialkan secara kontinu}{}{}.
\definitionsverweis {Grafik}{}{}
dari $f$ adalah

\mavergleichskettedisp
{\vergleichskette
{ Y
}
{ =} { \left\{ (x_1  , \ldots , x_{n-1},f( x_1  , \ldots , x_{n-1}))  \mid   x_1  , \ldots , x_{n-1} \in V         \right\}
}
{ \subseteq} { V \times \R
}
{ \subseteq} { \R^n
}
{ } {
}
}
{}{}{,}
yang juga dapat dipandang sebagai himpunan nol
\zusatzklammer {serat di atas $0$} {} {}
dari

\mavergleichskettedisp
{\vergleichskette
{ h(x_1  , \ldots , x_n)
}
{ \defeq} { x_n- f( x_1 , \ldots , x_{n-1})
}
{ } {
}
{ } {
}
{ } {
}
}
{}{}{.}
Turunan parsial dari $h$ adalah

\mathdisp {- { \frac{   \partial   f   }{  \partial   x_1   } } , \ldots , - { \frac{   \partial   f   }{  \partial   x_{n-1}   } } ,1} { . }

Khususnya, $h$
\definitionsverweis {reguler}{}{}
di setiap titik pada $Y$.
\definitionsverweis {Ruang tangen}{}{}
pada $Y$ di sebuah titik

\mavergleichskette
{\vergleichskette
{ P
}
{ \in }{ Y
}
{  }{
}
{    }{
}
{   }{
}
}
{}{}{}
tegak lurus terhadap
\mathl{\begin{pmatrix}  - { \frac{   \partial   f   }{  \partial   x_1   } } ( P)  \\ \vdots \\  -  { \frac{   \partial   f   }{  \partial   x_{n-1}   } } ( P)\\ 1   \end{pmatrix}}{.} Setiap garis dasar diparameterkan
\mathl{\begin{pmatrix}  x_1  \\ \vdots \\ x_{n-1}   \end{pmatrix} + t  \begin{pmatrix} v_1  \\\vdots\\ v_{n-1}   \end{pmatrix}}{} menjadi kurva yang diparameterkan
\maabbeledisp {} { I } { Y
} { t } { \begin{pmatrix}  x_1+tv_1  \\ \vdots \\  x_{n-1} +tv_{n-1}\\f \begin{pmatrix}  x_1+tv_1  \\ \vdots \\  x_{n-1} +tv_{n-1}   \end{pmatrix}   \end{pmatrix}
} {}
pada $Y$, yang turunannya menghasilkan sebuah vektor tangen. Jika lintasan untuk

\mavergleichskettedisp
{\vergleichskette
{ v
}
{ =} { e_i
}
{ } {
}
{ } {
}
{ } {
}
}
{}{}{}
dilambangkan dengan $\gamma_i$, maka

\mavergleichskettedisp
{\vergleichskette
{ \gamma'_i(t)
}
{ =} { \begin{pmatrix} \vdots \\ 1 \\ \vdots\\ { \frac{   \partial   f   }{  \partial   x_i   } }   \end{pmatrix}
}
{ } {
}
{ } {
}
{ } {
}
}
{}{}{}
dan

\mavergleichskettedisp
{\vergleichskette
{ \gamma^{\prime \prime}_i(t)
}
{ =} { \begin{pmatrix} \vdots \\ { \frac{   \partial    }{  \partial   x_i   } } { \frac{   \partial   f   }{  \partial   x_i   } }    \end{pmatrix}
}
{ } {
}
{ } {
}
{ } {
}
}
{}{}{}
menurut
[[Mehrere Variablen/k fach stetig differenzierbar/Längs einer Geraden/Fakt|Teorema 49.1 (Analisis (Osnabrück 2021--2023))]].

}

  \zwischenueberschrift{Analisis dan Geometri}

Misalkan

\mavergleichskette
{\vergleichskette
{ W
}
{ \subseteq }{ \R^n
}
{  }{
}
{    }{
}
{   }{
}
}
{}{}{}
\definitionsverweis {terbuka}{}{,}
\maabb {h} { W } { \R
} {}
adalah sebuah
\definitionsverweis {fungsi yang terdiferensialkan secara kontinu}{}{}
dan

\mavergleichskette
{\vergleichskette
{ Y
}
{ = }{ h^{-1}(c)
}
{  }{
}
{    }{
}
{   }{
}
}
{}{}{}
adalah
\definitionsverweis {serat}{}{}
di atas

\mavergleichskette
{\vergleichskette
{ c
}
{ \in }{ \R
}
{  }{
}
{    }{
}
{   }{
}
}
{}{}{,}
dengan $h$
\definitionsverweis {reguler}{}{}
di setiap titik pada $Y$. Kita membahas beberapa objek atau konstruksi yang memang ada di ruang linear sekitar $\R^n$, tetapi berhubungan sedemikian erat dengan hipermuka $Y$ sehingga kita ingin memahami dan mempelajarinya semata-mata sebagai objek pada $Y$. Untuk sementara, objek-objek ini memerlukan ruang sekitar karena definisinya mengacu pada analisis berdimensi lebih tinggi, yang saat ini baru tersedia untuk ruang vektor
\zusatzklammer {linear} {} {}.
Salah satu tujuan penting kita adalah mengembangkan konsep-konsep ini hanya pada objek geometris
\zusatzklammer {nonlinear} {} {}
$Y$. Kita akan menyinggung kurva pada hipermuka, ekstremum dengan kendala, dan persamaan diferensial untuk medan vektor tangensial.

\inputbemerkung
{}
{

Misalkan

\mavergleichskette
{\vergleichskette
{ W
}
{ \subseteq }{ \R^n
}
{  }{
}
{    }{
}
{   }{
}
}
{}{}{}
\definitionsverweis {terbuka}{}{,}
\maabb {h} { W } { \R
} {}
adalah sebuah
\definitionsverweis {fungsi yang terdiferensialkan secara kontinu}{}{}
dan

\mavergleichskette
{\vergleichskette
{ Y
}
{ = }{ h^{-1}(c)
}
{  }{
}
{    }{
}
{   }{
}
}
{}{}{}
adalah
\definitionsverweis {serat}{}{}
di atas

\mavergleichskette
{\vergleichskette
{ c
}
{ \in }{ \R
}
{  }{
}
{    }{
}
{   }{
}
}
{}{}{,}
dengan $h$
\definitionsverweis {reguler}{}{}
di setiap titik pada $Y$. Misalkan
\maabbdisp {\gamma} { I } { \R^n
} {}
adalah sebuah
\definitionsverweis {kurva terdiferensialkan}{}{,}
yang seluruhnya terletak di $Y$, dengan $I$ suatu interval terbuka. Maka, pada setiap saat

\mavergleichskette
{\vergleichskette
{ t
}
{ \in }{ I
}
{  }{
}
{    }{
}
{   }{
}
}
{}{}{}
\definitionsverweis {turunan}{}{}
$\gamma'(t)$ berada di ruang tangen pada $Y$ di titik $\gamma(t)$. Hal ini didasarkan pada sifat konstan dari

\mavergleichskette
{\vergleichskette
{ h \circ \gamma
}
{ = }{ c
}
{  }{
}
{    }{
}
{   }{
}
}
{}{}{,}
yang, melalui aturan rantai, memberikan

\mavergleichskettedisp
{\vergleichskette
{ \left(Dh \circ \gamma\right)_{t}
}
{ =} { \left(Dh  \right)_{ \gamma(t) }  \circ \left(D \gamma  \right)_{ t }
}
{ =} { \left(Dh  \right)_{ \gamma(t) } ( \gamma'( t) )
}
{ =} { 0
}
{ } {
}
}
{}{}{,}
dan karenanya

\mavergleichskette
{\vergleichskette
{ \gamma'( t)
}
{ \in }{ \operatorname{kern}  \left(Dh \right)_{ \gamma(t) }
}
{ = }{ T_{\gamma(t) }Y
}
{    }{
}
{   }{
}
}
{}{}{.}
Selain itu, dari
[[Tangentialraum/R/Faser/Kurvenrealisierung/Aufgabe|Soal 1.1]]
diperoleh bahwa setiap vektor tangen di $P$ pada $Y$ dapat direalisasikan oleh suatu kurva terdiferensialkan pada $Y$. Jadi, ruang tangen dapat dijelaskan secara bermakna hanya dengan kurva-kurva terdiferensialkan pada $Y$, meskipun keterdiferensialan kurva tersebut masih mengandaikan ruang sekitar.

}

\inputbemerkung
{}
{

Misalkan

\mavergleichskette
{\vergleichskette
{ W
}
{ \subseteq }{ \R^n
}
{  }{
}
{    }{
}
{   }{
}
}
{}{}{}
\definitionsverweis {terbuka}{}{,}
\maabb {h} { W } { \R
} {}
adalah sebuah
\definitionsverweis {fungsi yang terdiferensialkan secara kontinu}{}{}
dan

\mavergleichskette
{\vergleichskette
{ Y
}
{ = }{ h^{-1}(c)
}
{  }{
}
{    }{
}
{   }{
}
}
{}{}{}
adalah
\definitionsverweis {serat}{}{}
di atas

\mavergleichskette
{\vergleichskette
{ c
}
{ \in }{ \R
}
{  }{
}
{    }{
}
{   }{
}
}
{}{}{,}
dengan $h$
\definitionsverweis {reguler}{}{}
di setiap titik pada $Y$. Kita merangkum kembali teorema tentang ekstremum lokal dengan kendala dalam situasi ini; bandingkan
[[Extrema/Nebenbedingung/Hyperfläche/Fakt|Akibat 54.3 (Analisis (Osnabrück 2021--2023))]]
\zusatzklammer {dengan notasi yang berbeda} {} {.}
Selain fungsi $h$, yang menentukan $Y$ dan dengan demikian menentukan kendala, terdapat pula suatu lingkungan terbuka $U$ dari $Y$ beserta sebuah fungsi bernilai real $f$ yang didefinisikan di sana. Yang dicari adalah ekstremum lokal $f$, tetapi bukan pada seluruh $U$---yang biasanya diselidiki dengan gradien dan matriks Hesse---melainkan hanya pada $Y$. Misalnya, kita mungkin ingin mencari suhu maksimum pada permukaan suatu benda sambil mengabaikan bahwa bagian dalamnya lebih panas. Teorema tersebut menyatakan bahwa jika $f$ terdiferensialkan secara kontinu dan memiliki ekstremum lokal di suatu titik

\mavergleichskette
{\vergleichskette
{ P
}
{ \in }{ Y
}
{  }{
}
{    }{
}
{   }{
}
}
{}{}{}
di bawah kendala $Y$, maka gradien
\mathl{\operatorname{Grad} \,  f  (  P  )}{} merupakan kelipatan dari gradien
\mathl{\operatorname{Grad} \,  h  (  P  )}{}. Kriteria keterdiferensialan ini mengacu pada ruang sekitar: baik karena keterdiferensialan itu sendiri didefinisikan dengan mengacu pada ruang sekitar, maupun karena kedua gradien yang dibandingkan menjulur ke dalam ruang sekitar tersebut.

}

__NOEDITSECTION__

\inputfaktbeweis
{Gewöhnliches Differentialgleichungssystem/Hyperfläche/Lösung/Fakt}
{Teorema}
{}
{

\faktsituation {Misalkan

\mavergleichskette
{\vergleichskette
{ W
}
{ \subseteq }{ \R^n
}
{  }{
}
{    }{
}
{   }{
}
}
{}{}{}
\definitionsverweis {terbuka}{}{,}
\maabb {h} { W } { \R
} {}
adalah sebuah
\definitionsverweis {fungsi yang terdiferensialkan secara kontinu}{}{}
dan

\mavergleichskette
{\vergleichskette
{ Y
}
{ = }{ h^{-1}(c)
}
{  }{
}
{    }{
}
{   }{
}
}
{}{}{}
adalah
\definitionsverweis {serat}{}{}
di atas

\mavergleichskette
{\vergleichskette
{ c
}
{ \in }{ \R
}
{  }{
}
{    }{
}
{   }{
}
}
{}{}{,}
dengan $h$
\definitionsverweis {reguler}{}{}
di setiap titik pada $Y$.
Misalkan

\mavergleichskette
{\vergleichskette
{ Y
}
{ \subseteq }{ U
}
{ \subseteq }{ \R^n
}
{    }{
}
{   }{
}
}
{}{}{}
\definitionsverweis {terbuka}{}{}
dan misalkan $F$ adalah sebuah
\definitionsverweis {medan vektor}{}{}
terdiferensialkan pada $U$}
\faktvoraussetzung {dengan sifat bahwa untuk setiap titik

\mavergleichskette
{\vergleichskette
{ P
}
{ \in }{ Y
}
{  }{
}
{    }{
}
{   }{
}
}
{}{}{}
vektor $F(P)$ berada dalam
\definitionsverweis {ruang tangen pada serat}{}{}
di $P$ pada $Y$\zusatzfussnote {Medan vektor seperti ini disebut medan vektor tangensial} {.} {.}}
\faktfolgerung {Maka solusi $v(t)$ dari
\definitionsverweis {masalah nilai awal}{}{}
untuk $F$ dengan syarat awal

\mavergleichskette
{\vergleichskette
{ v(t_0)
}
{ \in }{ Y
}
{  }{
}
{    }{
}
{   }{
}
}
{}{}{}
seluruhnya terletak di $Y$.}
\faktzusatz {}
\faktzusatz {}

}
{

Kita bekerja dengan struktur Euklides pada $\R^n$ dan dengan
\definitionsverweis {gradien}{}{}
$\operatorname{Grad} \,  h$ dari $h$. Berdasarkan asumsi, pada $Y$ gradien ini tidak pernah sama dengan $0$, dan karena kekontinuan, sifat tersebut juga berlaku pada suatu lingkungan terbuka dari $Y$. Dengan memperkecil $U$ jika perlu, kita dapat menganggap bahwa gradien itu tidak memiliki titik nol di seluruh $U$. Pada $U$, kita tinjau medan vektor baru $G$ yang didefinisikan oleh

\mavergleichskettedisp
{\vergleichskette
{ G(P)
}
{ =} { F(P)-  { \frac{   \left\langle  F(P) ,
\operatorname{Grad} \,  h  (  P )    \right\rangle   }{  \Vert {
\operatorname{Grad} \,  h  (  P ) } \Vert^2   } }
\operatorname{Grad} \,  h  (  P )
}
{ } {
}
{ } {
}
{ } {
}
}
{}{}{.}
Untuk

\mavergleichskette
{\vergleichskette
{ P
}
{ \in }{ Y
}
{  }{
}
{    }{
}
{   }{
}
}
{}{}{}
berlaku

\mavergleichskette
{\vergleichskette
{ G(P)
}
{ = }{ F(P)
}
{  }{
}
{    }{
}
{   }{
}
}
{}{}{,}
karena pada $Y$ gradien $h$ tegak lurus terhadap medan vektor tersebut. Selain itu, $G$
\zusatzklammer {berbeda dengan $F$} {} {}
memiliki sifat bahwa untuk semua titik

\mavergleichskette
{\vergleichskette
{ P
}
{ \in }{ U
}
{  }{
}
{    }{
}
{   }{
}
}
{}{}{}
vektor $G(P)$ tegak lurus terhadap gradien, sebab

\mavergleichskettealignhandlinks
{\vergleichskettealignhandlinks
{ \left\langle  G(P) ,
\operatorname{Grad} \,  h  (  P )   \right\rangle
}
{ =} { \left\langle  F(P)- { \frac{   \left\langle  F(P) ,
\operatorname{Grad} \,  h  (  P )    \right\rangle  }{  \Vert {
\operatorname{Grad} \,  h  (  P ) } \Vert^2  } }
\operatorname{Grad} \,  h  (  P )  ,
\operatorname{Grad} \,  h  (  P )   \right\rangle
}
{ =} { \left\langle  F(P) ,
\operatorname{Grad} \,  h  (  P )   \right\rangle - \left\langle  F(P) ,
\operatorname{Grad} \,  h  (  P )   \right\rangle
}
{ =} { 0
}
{ } {
}
}
{}
{}{.}
Sekarang misalkan
\maabbdisp {v} { I} { \R^n
} {}
adalah, menurut
[[Picard Lindelöf/Lokal Lipschitz/Lokale Existenz und Eindeutigkeit/Fakt|Teorema 56.2 (Analisis (Osnabrück 2021--2023))]]
\zusatzfussnote {Catatan edisi: Dengan asumsi yang dinyatakan, fungsi h hanya berkelas C¹ sehingga gradiennya sekadar kontinu; hal ini belum membuktikan bahwa medan G lokal Lipschitz sebagaimana diperlukan untuk kesimpulan keunikan Picard--Lindelöf. Langkah ini memerlukan regularitas tambahan atau pembuktian lokal melalui koordinat manifold.} {} {}
,
solusi tunggal dari masalah nilai awal untuk $G$ dengan syarat awal
\mavergleichskette
{\vergleichskette
{ v(t_0)
}
{ = }{ P
}
{ \in }{ Y
}
{    }{
}
{   }{
}
}
{}{}{.}
Maka

\mavergleichskettealign
{\vergleichskettealign
{ \left(Dh \circ v \right)_{t}
}
{ =} { \left(Dh  \right)_{v(t)} \circ  \left(D v \right)_{t}
}
{ =} { \left(Dh  \right)_{v(t)}  \left(  v'(t)  \right)
}
{ =} { \left(Dh  \right)_{v(t)}  \left(  G(v(t))  \right)
}
{ =} { \left\langle
\operatorname{Grad} \,  h  ( v(t) )  ,  G(v(t))    \right\rangle
}
}
{
\vergleichskettefortsetzungalign
{ =} { 0
}
{ } {}
{ } {}
{ } {}
}
{}{.}
Oleh karena itu, $h$ konstan pada citra $v( I)$, dan karena

\mavergleichskette
{\vergleichskette
{ h(v(t_0))
}
{ = }{ c
}
{  }{
}
{    }{
}
{   }{
}
}
{}{}{}
maka

\mavergleichskette
{\vergleichskette
{ h(v(t))
}
{ = }{ c
}
{  }{
}
{    }{
}
{   }{
}
}
{}{}{}
untuk semua

\mavergleichskette
{\vergleichskette
{ t
}
{ \in }{ I
}
{  }{
}
{    }{
}
{   }{
}
}
{}{}{,}
sehingga

\mavergleichskette
{\vergleichskette
{ v(t)
}
{ \in }{ Y
}
{  }{
}
{    }{
}
{   }{
}
}
{}{}{}
untuk semua

\mavergleichskette
{\vergleichskette
{ t
}
{ \in }{ I
}
{  }{
}
{    }{
}
{   }{
}
}
{}{}{.}

}

Dengan regularitas tambahan, cara lain untuk membuktikan pernyataan ini adalah dengan memindahkan medan vektor $F$ melalui suatu difeomorfisme lokal antara

\mavergleichskette
{\vergleichskette
{ P
}
{ \in }{ V
}
{ \subseteq }{ Y
}
{    }{
}
{   }{
}
}
{}{}{}
dan

\mavergleichskette
{\vergleichskette
{ V'
}
{ \subseteq }{ \R^{n-1}
}
{  }{
}
{    }{
}
{   }{
}
}
{}{}{}
ke $\R^{n-1}$, memastikan bahwa medan dalam koordinat bersifat lokal Lipschitz, membuktikan keberadaan solusi di sana, memindahkan kembali solusi itu ke $Y$, lalu merekatkan solusi-solusi parsial tersebut. Tanpa regularitas tersebut, argumen alternatif ini belum lengkap.

Lingkungan terbuka $U$ dari $Y$ diperlukan agar kita dapat berbicara tentang medan vektor terdiferensialkan, meskipun pada akhirnya hanya medan vektor pada $Y$ yang relevan.

  \zwischenueberschrift{Percepatan}

\inputbemerkung
{}
{

Kita tinjau kurva
\maabbeledisp {v} {\R} { \R^2
} { t } { \begin{pmatrix}  \cos  t  \\ \sin  t    \end{pmatrix}
} {,}
yang seluruhnya bergerak pada lingkaran satuan. Dengan demikian, turunannya

\mavergleichskette
{\vergleichskette
{ v'(t)
}
{ = }{ \begin{pmatrix}  - \sin  t  \\ \cos  t    \end{pmatrix}
}
{  }{
}
{    }{
}
{   }{
}
}
{}{}{}
selalu tangensial terhadap lingkaran satuan
\zusatzklammer {hal ini juga mengikuti dari
[[Reguläre Hyperfläche/Kurven/Tangentiale Realisierung/Motivation/Bemerkung|Catatan 1.1.3]]} {} {.}
Normanya selalu sama dengan $1$, sehingga lingkaran satuan ditempuh secara seragam dengan besar kecepatan konstan. Namun, arah kecepatannya berubah secara kontinu, dan turunan keduanya adalah

\mavergleichskettedisp
{\vergleichskette
{ v^{\prime \prime}(t)
}
{ =} { \begin{pmatrix}  - \cos  t  \\ - \sin  t    \end{pmatrix}
}
{ =} { - \begin{pmatrix}  \cos  t  \\ \sin  t    \end{pmatrix}
}
{ =} { - v(t)
}
{ } {
}
}
{}{}{.}
Percepatannya adalah negatif dari vektor posisi dan tegak lurus terhadap vektor kecepatan. Jika ruang tangen dari bidang sekitar di sebuah titik $P$ pada lingkaran---yakni $\R^2$ yang dipandang berpusat di $P$---diuraikan menjadi arah yang tangensial terhadap lingkaran
\zusatzklammer {ruang tangen pada serat} {} {}
dan arah yang tegak lurus terhadapnya
\zusatzklammer {ruang normal} {} {}
\zusatzklammer {dalam arti penguraian suatu ruang vektor sebagai
\definitionsverweis {jumlah langsung}{}{}
dari subruang-subruang vektor} {} {,}
maka seluruh percepatan berlangsung di ruang normal; proyeksi ortogonalnya pada ruang tangen selalu $0$. Dalam pengertian ini, percepatan tidak relevan bagi gerak pada lingkaran satuan.

Sekarang kita tinjau secara lebih umum sebuah kurva
\maabbeledisp {w} {\R} { \R^2
} { t } { \begin{pmatrix}  \cos  f(t)  \\ \sin f(t)    \end{pmatrix}
} {,}
dengan $f$ suatu fungsi real yang terdiferensialkan dua kali. Seperti sebelumnya, kurva ini bergerak pada lingkaran satuan, dan turunannya adalah

\mavergleichskettedisp
{\vergleichskette
{ w'(t)
}
{ =} { f'(t) \begin{pmatrix}  - \sin  f(t)  \\ \cos f(t)    \end{pmatrix}
}
{ } {
}
{ } {
}
{ } {
}
}
{}{}{,}
sehingga, seperti sebelumnya, kurva ini selalu tangensial terhadap lingkaran satuan, sedangkan norma kecepatannya sekarang sama dengan $\betrag { f'(t) }$. Turunan keduanya adalah

\mavergleichskettedisp
{\vergleichskette
{ w^{\prime \prime} (t)
}
{ =} { f^{\prime \prime} (t) \begin{pmatrix}  - \sin  f(t)  \\ \cos f(t)    \end{pmatrix} - \left(  f'(t)  \right)^2 \begin{pmatrix}  \cos  f(t)  \\ \sin f(t)    \end{pmatrix}
}
{ } {
}
{ } {
}
{ } {
}
}
{}{}{,}
yang secara langsung memperlihatkan penguraian ke dalam komponen tangensial dan komponen yang ortogonal terhadapnya. Suku pertama, yang memuat turunan kedua fungsi, adalah komponen tangensial. Komponen normalnya adalah
\mathl{- \left(  f'(t)  \right)^2 \begin{pmatrix}  \cos  f(t)  \\ \sin f(t)    \end{pmatrix}}{,} sedangkan percepatan tangensial muncul ketika turunan kedua $f$ tidak sama dengan nol.

}

Untuk sebuah hipermuka

\mavergleichskette
{\vergleichskette
{ Y
}
{ \subseteq }{ \R^n
}
{  }{
}
{    }{
}
{   }{
}
}
{}{}{}
dan sebuah titik

\mavergleichskette
{\vergleichskette
{ P
}
{ \in }{ Y
}
{  }{
}
{    }{
}
{   }{
}
}
{}{}{}
kita selalu dapat menguraikan ruang sekitar $\R^n$, yang dipandang sebagai ruang tangen dari $\R^n$ di $P$, secara ortogonal sebagai

\mavergleichskettedisp
{\vergleichskette
{ \R^n
}
{ =} { T_PY \oplus N_PY
}
{ } {
}
{ } {
}
{ } {
}
}
{}{}{,}
dengan $N_PY$ sebuah garis yang terdiri atas semua titik yang ortogonal terhadap $T_PY$ dan disebut \stichwort {garis normal} {}. Garis normal terdiri atas semua kelipatan gradien
\mathl{\operatorname{Grad} \,  h  (  P )}{} di $P$, jika $h$ menyatakan suatu fungsi yang mendeskripsikan $Y$. Sebuah \stichwort {vektor normal} {} adalah elemen garis normal bernorma $1$. Ada dua vektor normal demikian, dan yang satu merupakan negatif dari yang lain. Melalui
\definitionsverweis {proyeksi ortogonal}{}{}
sepanjang $N_PY$, yakni pemetaan
\maabbdisp {\pi_P} {\R^n } { T_PY
} {}
kita dapat memetakan setiap vektor di $\R^n$ ke sebuah vektor di ruang tangen.

\inputdefinition
{}
{

Misalkan

\mavergleichskette
{\vergleichskette
{ W
}
{ \subseteq }{ \R^n
}
{  }{
}
{    }{
}
{   }{
}
}
{}{}{}
\definitionsverweis {terbuka}{}{,}
\maabb {h} { W } { \R
} {}
adalah sebuah
\definitionsverweis {fungsi yang terdiferensialkan secara kontinu}{}{}
dan

\mavergleichskette
{\vergleichskette
{ Y
}
{ = }{ h^{-1}(c)
}
{  }{
}
{    }{
}
{   }{
}
}
{}{}{}
adalah
\definitionsverweis {serat}{}{}
di atas

\mavergleichskette
{\vergleichskette
{ c
}
{ \in }{ \R
}
{  }{
}
{    }{
}
{   }{
}
}
{}{}{,}
dengan $h$
\definitionsverweis {reguler}{}{}
di setiap titik pada $Y$. Misalkan
\maabbdisp {\gamma} {I} {Y
} {}
merupakan
\zusatzklammer {sebagai pemetaan ke $\R^n$} {} {}
sebuah \definitionsverweis {kurva yang terdiferensialkan dua kali secara kontinu}{}{.}
Maka pemetaan
\maabbeledisp {} {I} { TY
} {t} { \pi_{\gamma(t)}( \gamma^{\prime \prime} (t))
} {,}
dengan $\pi_{\gamma(t)}$ menyatakan proyeksi ortogonal
\maabbdisp {} {\R^n } { T_{\gamma(t)}Y
} {}
disebut
\definitionswort {turunan tangensial kedua}{}
\zusatzklammer {atau
\definitionswort {percepatan tangensial}{}} {} {}
dari $\gamma$.

}

Gagasan tentang turunan kedua yang tangensial terhadap hipermuka dikembangkan lebih lanjut dalam konteks yang disebut koneksi.

  \zwischenueberschrift{Kurva Geodesik}

\inputdefinition
{}
{

Misalkan

\mavergleichskette
{\vergleichskette
{ W
}
{ \subseteq }{ \R^n
}
{  }{
}
{    }{
}
{   }{
}
}
{}{}{}
\definitionsverweis {terbuka}{}{,}
\maabb {h} { W } { \R
} {}
adalah sebuah
\definitionsverweis {fungsi yang terdiferensialkan secara kontinu}{}{}
dan

\mavergleichskette
{\vergleichskette
{ Y
}
{ = }{ h^{-1}(c)
}
{  }{
}
{    }{
}
{   }{
}
}
{}{}{}
adalah
\definitionsverweis {serat}{}{}
di atas

\mavergleichskette
{\vergleichskette
{ c
}
{ \in }{ \R
}
{  }{
}
{    }{
}
{   }{
}
}
{}{}{,}
dengan $h$
\definitionsverweis {reguler}{}{}
di setiap titik pada $Y$. Sebuah
\zusatzklammer {sebagai pemetaan ke $\R^n$} {} {}
\definitionsverweis {kurva yang terdiferensialkan dua kali secara kontinu}{}{}
\maabbdisp {\gamma} {I} {Y
} {}
disebut
\definitionswort {geodesik}{}
\zusatzklammer {atau
\definitionswort {kurva geodesik}{}
pada $Y$} {} {,}
jika
\definitionsverweis {percepatan tangensial}{}{}
kurva itu sama dengan $0$ di mana-mana.

}

Sebuah kurva yang terdiferensialkan dua kali merupakan geodesik tepat ketika turunan keduanya selalu ortogonal terhadap ruang tangen. Gagasan fisika yang mendasarinya adalah gerak sebuah partikel yang bergerak pada $Y$ tanpa percepatan. Secara umum memang terdapat gaya-gaya kendala, seperti gravitasi atau tarikan magnetis, yang memaksa geraknya tetap pada $Y$ dan pada gilirannya memerlukan percepatan di ruang sekitar; namun pada $Y$ sendiri tidak ada percepatan. Jadi, kurva geodesik menggambarkan gerak alami tanpa percepatan pada sebuah hipermuka terdiferensialkan.

\bild{ \begin{center}
\includegraphics[width=5.5cm]{ \bildeinlesungsvg {Great circle passing through two points} {svg} }
\end{center}
\bildtext {} }

\bildlizenz { Great circle passing through two points.svg } {} {HaEr48} {Commons} {CC BY-SA 4.0} {}

\inputdefinition
{}
{

Irisan sebuah
\definitionsverweis {permukaan bola}{}{}
\mathl{S}{} dengan suatu bidang yang melalui pusat bola disebut sebuah \definitionswort {lingkaran besar}{} pada $S$.

}

Pada bola Bumi, khatulistiwa dan lingkaran-lingkaran bujur merupakan lingkaran besar; lingkaran-lingkaran lintang pada umumnya bukan.

\inputbeispiel{}
{

Kita berada pada bola satuan

\mavergleichskettedisp
{\vergleichskette
{ Y
}
{ =} { \left\{ (x,y,z)  \mid   x^2+y^2+z^2 = 1         \right\}
}
{ \subseteq} { \R^3
}
{ } {
}
{ } {
}
}
{}{}{.}
Di

\mavergleichskette
{\vergleichskette
{ P
}
{ \in }{ Y
}
{  }{
}
{    }{
}
{   }{
}
}
{}{}{}
terdapat sebuah partikel yang menerima impuls gerak dalam arah tangensial tertentu yang dinyatakan oleh

\mavergleichskette
{\vergleichskette
{ v
}
{ \in }{ T_PY
}
{  }{
}
{    }{
}
{   }{
}
}
{}{}{.}
Partikel itu terus bergerak dalam arah ini
\zusatzklammer {tanpa diperlambat dan tanpa dipercepat} {} {}
namun harus selalu tetap berada pada bola
\zusatzklammer {karena gravitasi Bumi atau karena permukaan yang kukuh} {} {.}
Misalkan titik tersebut diberikan oleh
\mathl{\begin{pmatrix}  x_1  \\ y_1 \\ z_1   \end{pmatrix}}{} dan vektor tangen tegak lurus yang menyatakan impuls sesaat diberikan oleh

\mavergleichskettedisp
{\vergleichskette
{ \begin{pmatrix}  x_3  \\ y_3 \\ z_3   \end{pmatrix}
}
{ =} { a \begin{pmatrix}  x_2  \\ y_2 \\ z_2   \end{pmatrix}
}
{ } {
}
{ } {
}
{ } {
}
}
{}{}{,}
dengan $\begin{pmatrix}  x_2  \\ y_2 \\ z_2   \end{pmatrix}$ suatu vektor bernorma $1$ yang tegak lurus terhadap vektor posisi. Gerak tersebut berada pada bidang yang ditentukan oleh kedua vektor ini sekaligus pada bola. Salah satu parametrisasi geraknya adalah
\maabbeledisp {\gamma} {\R} { \R^3
} { t } { \cos \left( at \right) \begin{pmatrix}  x_1  \\ y_1 \\ z_1   \end{pmatrix} + \sin  \left( at \right) \begin{pmatrix}  x_2  \\ y_2 \\ z_2   \end{pmatrix}
} {.}
Berlaku

\mavergleichskette
{\vergleichskette
{ \gamma(0)
}
{ = }{ \begin{pmatrix}  x_1  \\ y_1 \\ z_1   \end{pmatrix}
}
{  }{
}
{    }{
}
{   }{
}
}
{}{}{}
dan

\mavergleichskette
{\vergleichskette
{ \gamma'(t)
}
{ = }{ -a \sin  \left( at \right) \begin{pmatrix}  x_1  \\ y_1 \\ z_1   \end{pmatrix} +a \cos \left( at \right) \begin{pmatrix}  x_2  \\ y_2 \\ z_2   \end{pmatrix}
}
{  }{
}
{    }{
}
{   }{
}
}
{}{}{,}
sehingga

\mavergleichskette
{\vergleichskette
{ \gamma'(0)
}
{ = }{ a \begin{pmatrix}  x_2  \\ y_2 \\ z_2   \end{pmatrix}
}
{  }{
}
{    }{
}
{   }{
}
}
{}{}{.}

}

\inputbeispiel{}
{

Kita tinjau selimut silinder

\mavergleichskettedisp
{\vergleichskette
{ Z
}
{ =} { \left\{  \begin{pmatrix}  x  \\ y \\ z   \end{pmatrix} \in \R^3  \mid   x^2+y^2 = 1         \right\}
}
{ } {
}
{ } {
}
{ } {
}
}
{}{}{.}
Secara geometris, kita memperoleh gerakan-gerakan berikut pada silinder, yang percepatannya selalu ortogonal terhadap ruang tangen. Ruang tangen di titik
\mathl{\begin{pmatrix}  x  \\ y \\ z   \end{pmatrix}}{} diberikan oleh
\mathl{\R  \begin{pmatrix}  y  \\ -x\\  0   \end{pmatrix} + \R \begin{pmatrix}  0  \\ 0\\  1   \end{pmatrix}}{}, dan vektor $- \begin{pmatrix}  x  \\ y \\  0   \end{pmatrix}$ tegak lurus terhadapnya. Untuk gerak-gerak berikut, penskalaan parameter waktunya dapat diubah secara afin.
\aufzaehlungdrei{Gerak pada sebuah garis di selimut silinder yang sejajar dengan sumbu silinder:

\mavergleichskettedisp
{\vergleichskette
{ \gamma(t)
}
{ =} { \begin{pmatrix}  x_0  \\ y_0 \\ t   \end{pmatrix}
}
{ } {
}
{ } {
}
{ } {
}
}
{}{}{}
dengan

\mavergleichskette
{\vergleichskette
{ x_0^2 +y_0^2
}
{ = }{ 1
}
{  }{
}
{    }{
}
{   }{
}
}
{}{}{.}
Turunan pertamanya adalah $\begin{pmatrix}  0 \\ 0\\  1   \end{pmatrix}$ dan turunan keduanya adalah $0$.
}{Gerak melingkar yang diberikan oleh

\mavergleichskettedisp
{\vergleichskette
{ \delta(t)
}
{ =} { \begin{pmatrix}  \cos  t  \\ \sin  t \\ c   \end{pmatrix}
}
{ } {
}
{ } {
}
{ } {
}
}
{}{}{.}
Turunan pertamanya adalah $\begin{pmatrix}  - \sin  t  \\ \cos  t \\  0   \end{pmatrix}$ dan turunan keduanya adalah $- \begin{pmatrix}  \cos  t  \\ \sin  t \\  0   \end{pmatrix}$.
}{

\bild{ \begin{center}
\includegraphics[width=5.5cm]{ \bildeinlesungjpg { 2019-07-Helix} {jpg} }
\end{center}
\bildtext {} }

\bildlizenz {  2019-07-Helix.jpg } {} {Ag2gaeh} {Commons} {CC BY-SA 4.0} {}

Sebuah kurva ulir
\zusatzklammer {heliks} {} {}
yang diberikan oleh

\mavergleichskettedisp
{\vergleichskette
{ \epsilon(t)
}
{ =} { \begin{pmatrix}  \cos  t  \\ \sin  t \\  e t   \end{pmatrix}
}
{ } {
}
{ } {
}
{ } {
}
}
{}{}{}
dengan

\mavergleichskette
{\vergleichskette
{ e
}
{ \neq }{ 0
}
{  }{
}
{    }{
}
{   }{
}
}
{}{}{.}
Turunan pertamanya adalah $\begin{pmatrix}  - \sin  t  \\ \cos  t \\ e   \end{pmatrix}$ dan turunan keduanya adalah $- \begin{pmatrix}  \cos  t  \\ \sin  t \\  0   \end{pmatrix}$.
}

}

 [[Kategorie:Latexseite]]
